\section{The Four Layers Are Not a Ranking}

\textbf{[STIPULATION]} The entire series marks every statement with one of four
layers: \textbf{[K]} core (derived from $\xi$), \textbf{[B]} bridge (algebraically
proved), \textbf{[S]} sketch (plausibility), \textbf{[STIPULATION]} declared. This
document says what this marking means and what it does not mean.

It is an \emph{ordering}, not a \emph{ranking}. The layers say with what right a
statement stands -- not how valuable it is. A \textbf{[STIPULATION]} is not
inferior to a \textbf{[K]}; it is of a different kind. Without stipulations there
would be no core from which anything could follow.

\textbf{[B]} The confusion of ordering with ranking is the most common
misreading of layered presentations. One who reads \textbf{[S]} as ``not yet quite
[K]'' assumes that everything strives toward the core. It does not. A sketch
marked as a sketch is complete in its place -- it claims no more than it shows,
and that is its strength.

\section{What Each Layer Accomplishes}

\textbf{[STIPULATION]} The four layers in their own language:

\textbf{[K] Core.} From the three stipulations (A020) the statement follows by
computation or geometric necessity. Examples: the non-closure theorem (A050),
$Q=2/3$ from $r=\sqrt2$ (A110), $\Delta w = 1$ (A180). The core is what the
theory stands or falls by.

\textbf{[B] Bridge.} The statement is algebraically proved, but uses a
translation or a known structure that does not itself follow from $\xi$. Example:
the Hilbert space representation (A070). A bridge carries everything on both its
sides and proves nothing new -- it connects.

\textbf{[S] Sketch.} The statement is plausible but not computed through.
Examples: the Landauer connection (A180), the path assignment of entanglement
(A160). A sketch is an honest placeholder, not a half-truth.

\textbf{[STIPULATION] Declaration.} The statement is stipulated, not derived.
Examples: the three fundamental stipulations (A020), the choice of the
characteristic mass (A145). A stipulation is the entry point without which a
closed system cannot begin (A130).

\section{The Self-Application}

\textbf{[STIPULATION]} The ordering applies to itself. The fundamental relation
$\tilde T\cdot m = 1$ is a \textbf{[STIPULATION]}, not a \textbf{[K]}. It is not
proved; it is the beginning.

This is not a weakness but the condition of possibility of the whole. Every theory
that derives something derives it from something it does not derive. The difference
between a clean and an unclean theory is not whether it has stipulations -- every
theory does -- but whether it \emph{declares them as such}.

\textbf{[B]} That is exactly what the layer marking does: it makes the stipulations
visible instead of disguising them as results. The error that A130 identifies in
the older $\alpha$ presentation -- presenting a calibration as a derivation -- is a
violation of precisely this rule. The marking is the means to avoid it.

\section{Why No Ranking Is Possible}

\textbf{[B]} One might try to rank the layers: [K] better than [B] better than [S]
better than [STIPULATION]. The attempt fails at the self-application. The most
valuable statement of the theory -- $\tilde T\cdot m = 1$ -- would then be the
least. That is absurd, and the absurdity shows that the ordering is horizontal, not
vertical.

The layers are four ways of being true, not four degrees of truth.

\section{Open Questions}

\textbf{First, the assignment in individual cases is not always sharp.} Some
statements lie between [B] and [S]. The series decides such cases conservatively
-- when in doubt, the weaker layer -- but the boundary is not compelling in every
case.

\textbf{Second, the marking does not replace judgment.} It says with what right a
statement stands; whether this right is correctly assigned in the individual case
must be checked. The verification scripts do this for the computable cases.

\section{Supersedes}

This document takes up the methodological note ``Ordnung ohne Rangordnung / Order
without Ranking'' and makes its distinction (law, norm, stipulation) the
throughgoing principle of the series. It stands alongside A010, which names the
rules; A200 justifies them.

\section{Use of AI Tools}

This document was created with the assistance of AI tools. The questions,
stipulations, and all substantive decisions come from the author; the tools
formulate, compute, and create verification scripts.
Responsibility for content and errors rests entirely with the author.
Full wording of the rules in A010.
